\chapter{Exit codes}\label{app:exit-codes}
{\small
\begin{longtable}{@{}>{\raggedright\arraybackslash}p{0.24\linewidth}>{\raggedright\arraybackslash}p{0.72\linewidth}@{}}
\toprule
\textbf{Command} & \textbf{Exit codes} \\ \midrule\endhead
all kit commands & 0 success; 2 usage error (unknown command or option, invalid input) \\
\texttt{test} & 0 all checks passed; otherwise the exit code of the first failing step \\
\texttt{simulate} & 0 all expectations met; 1 at least one mismatch; 2 model or expectations error \\
\texttt{verify} & 0 all checked expectations met; 1 at least one mismatch; 2 a CLI call failed \\
\texttt{export} & 0 written (read \texttt{REPORT.md}); 2 a failed read, an HTTP error or invalid options \\
\texttt{forwarder} & 0 written; 2 invalid combination or duplicate name (nothing is written) \\
\texttt{conjur} & 0 written; 2 unreadable or invalid policy file \\
\texttt{account-settings} & 2 invalid settings; otherwise the exit code of \texttt{akeyless update-account-settings}; 0 with \texttt{--dry-run} \\
\texttt{token} & 0 token printed; 1 missing settings or authentication failed \\
\texttt{build-skill} & 0 packages written; 2 the built skill failed validation \\
\texttt{install-skill} & 0 installed (or printed with \texttt{--dry-run}); 2 target exists without \texttt{--force}, or \texttt{custom} without \texttt{--dest} \\
\texttt{validate-skill} & 0 valid; 1 problems found (listed) \\
\texttt{mcp} & runs until the client closes stdin \\
CI drift job & \texttt{plan -detailed-exitcode}: 0 no changes, 2 drift (job fails), 1 error \\
\bottomrule
\end{longtable}}

\chapter{Environment variables}\label{app:env}
{\small
\begin{longtable}{@{}>{\raggedright\arraybackslash}p{0.34\linewidth}>{\raggedright\arraybackslash}p{0.62\linewidth}@{}}
\toprule
\textbf{Variable} & \textbf{Used by and meaning} \\ \midrule\endhead
\texttt{AKEYLESS\_API} & \texttt{token}, \texttt{export --transport api}: API base URL (default \texttt{https://api.akeyless.io}, or a Gateway \texttt{/api/v2}) \\
\texttt{AKEYLESS\_ACCESS\_ID} & \texttt{token}: access ID of the CI auth method. Provider and \texttt{export --hcl --generate}: with \texttt{AKEYLESS\_ACCESS\_KEY}, API-key login \\
\texttt{AKEYLESS\_ACCESS\_KEY} & provider API-key login (bootstrap, \texttt{--generate}); never in CI for the main pipeline \\
\texttt{CI\_JWT} & \texttt{token}: the CI-issued JWT \\
\texttt{AKEYLESS\_TOKEN} & \texttt{export --transport api}, \texttt{account-settings}: Akeyless token when \texttt{--token} is not given \\
\texttt{AKEYLESS\_GATEWAY} & provider fallback for the API endpoint \\
\texttt{TF\_VAR\_pipeline\_jwt} & Terraform: the CI JWT (must be set; see troubleshooting) \\
\texttt{TF\_VAR\_pipeline\_access\_id} & Terraform: the CI auth method's access ID (CI variable \texttt{AKEYLESS\_PIPELINE\_ACCESS\_ID}) \\
\texttt{TF\_VAR\_auth\_method\_secrets} & JSON map for \texttt{*\_var} keys in \texttt{auth\_methods.yaml} \\
\texttt{TF\_VAR\_event\_forwarder\_secrets} & JSON map for \texttt{*\_var} keys of event forwarders \\
\texttt{TF\_VAR\_gateway\_secrets} & JSON map for \texttt{*\_var} keys in \texttt{gateway.yaml} \\
\texttt{TF\_VAR\_gateway\_api} & Gateway V2 API for gateway objects (CI variable \texttt{AKEYLESS\_GATEWAY\_API}) \\
\texttt{WATCHOVER\_REPO} & \texttt{mcp}: kit root the tools run in (default: the repository of the server) \\
\texttt{WATCHOVER\_MCP\_LIVE} & \texttt{mcp}: \texttt{1} also exposes the commands that call Akeyless \\
\bottomrule
\end{longtable}}

\chapter{Repository layout}\label{app:layout}
\begin{lstlisting}
terraform/            data-driven Terraform (provider akeyless-community/akeyless ~> 2.0)
  versions.tf         provider pin, JWT login, gateway alias, endpoint variables
  main.tf             inputs, model loading, guards, folders
  auth_methods.tf roles.tf governance.tf pki.tf gateway.tf outputs.tf
  model/              THE MODEL: roles/<bu>.yaml, auth_methods, folders, governance, pki,
                      gateway, access_expectations
  tests/              terraform test (mocked provider), fixtures/, guard_messages.sh,
                      export_roundtrip.sh
bootstrap/            one-time CI identity (OAuth2/JWT) + plan/apply roles, own state
scripts/              watchover (entry point) and the tools behind it
  tests/              pytest suite, fake CLI/REST back ends, fixtures
ci/                   GitHub Actions and GitLab CI templates, CODEOWNERS example
skill/                skill source: SKILL.md (Agent Skills format), agents/openai.yaml sidecar
AGENTS.md             instructions for AI coding agents working on this repository
                      (CLAUDE.md and GEMINI.md import it)
docs/                 Markdown documentation (bundled in the skill), this PDF
  manual/             LaTeX sources of this manual, build.py, demo back end
\end{lstlisting}

\chapter{Glossary}\label{app:glossary}
\begin{description}[style=nextline,leftmargin=1.5em,font=\normalfont\bfseries]
\item[Agent Skills] Open format for packaging instructions, references and scripts as a folder with a \texttt{SKILL.md}; read by many AI agents (\url{https://agentskills.io}).
\item[Access role] An Akeyless role: a set of rules (path + capabilities) plus management-plane permissions.
\item[Association] The binding of a role to an auth method, optionally narrowed by sub-claims.
\item[Auth method] How an identity authenticates (OIDC, SAML, AWS IAM, Kubernetes, API key, \ldots); it carries an access ID.
\item[Anti-escalation] Akeyless refuses to let an identity grant rights it does not hold itself.
\item[Authoritative] Terraform owns the complete list: rules not in the model are removed on apply.
\item[Capability] An allowed action on a path: read, list, create, update, delete, deny, or SRA capabilities.
\item[Drift] A difference between the account and Git, detected by a non-empty plan.
\item[ForceNew] A provider attribute whose change replaces the object (destroy + create).
\item[Guard] A Terraform precondition that rejects an invalid or unsafe model before any change.
\item[Management plane] Who may administer Akeyless itself (audit, analytics, gateways, event centre \ldots), as opposed to accessing secrets.
\item[MCP] Model Context Protocol: lets an AI client use a server's tools, resources and prompts. \texttt{watchover mcp} is such a server.
\item[Model] The YAML files under \texttt{terraform/model/}: the single source of truth.
\item[Persona] A test identity (auth method + claims) used in access expectations.
\item[Prompt pack] The skill as plain instructions and knowledge files, for assistants without skill or MCP support.
\item[Rule type] What a rule applies to: items, targets, roles, auth methods, SRA, ARA.
\item[SRA] Secure Remote Access: brokered sessions to infrastructure.
\item[Sub-claims] Claim values an identity must present for an association to apply (OR within a key, AND across keys).
\item[VERIFY] A statement not confirmed by a primary source; check it on your account.
\end{description}

\chapter{Sources}\label{app:sources}
\begin{itemize}
\item Akeyless documentation: \url{https://docs.akeyless.io/}
\item Akeyless OpenAPI specification and Go SDK v5.0.x (the API operations and their request schemas)
\item Terraform provider \texttt{akeyless-community/akeyless} v2.0.5: source code and schema
  (\url{https://registry.terraform.io/providers/akeyless-community/akeyless/latest/docs})
\item Agent Skills specification: \url{https://agentskills.io/specification}; the skills and MCP documentation of
  the agents listed in Chapter~\ref{doc-agents}; Model Context Protocol: \url{https://modelcontextprotocol.io}
\item Terraform and OpenTofu documentation for \texttt{terraform test}, mocked providers and import blocks
\end{itemize}
Statements that no primary source confirms are marked \textbf{VERIFY} (Chapter~\ref{doc-limitations}).
